\title{Byte Latent Transformer: Patches Scale Better Than Tokens}

\begin{abstract}
We present the Byte Latent Transformer (\textsc{BLT}), a novel large language model architecture that operates at the byte level and achieves, for the first time, comparable performance to tokenization-dependent LLMs at scale, while offering notable gains in both inference efficiency and robustness. The \textsc{BLT} framework processes input by transforming bytes into variably sized patches, which constitute the model's fundamental computational units. These patches are delineated according to the entropy of the upcoming byte, dedicating greater computational resources and model expressivity to more complex, information-rich sequences. We conduct the first computational budget-aware scaling investigation for byte-level architectures, reaching up to 8B parameters and 4T training bytes. The findings confirm that it is practical to train models directly on byte-level data without relying on a predetermined vocabulary. Efficiency is enhanced in both training and inference stages as the model adaptively uses longer patches in contexts with high predictability, and our approach also delivers qualitative advances in reasoning ability and long tail generalization. In summary, at equivalent inference cost, \textsc{BLT} surpasses the scaling capabilities of token-based systems by simultaneously expanding both patch length and model parameters.
\end{abstract}

\section{Introduction}
\label{section:intro}

We present the Byte Latent Transformer~(\textbf{BLT}), a novel, tokenizer-independent framework designed to operate directly on raw byte sequences. For the first time, BLT attains comparable large-scale performance to traditional models that rely on tokenization, while offering notable gains in computational efficiency and resilience (\S\ref{sec:robustness}). In contrast to current large language models (LLMs), which are trained end-to-end except for the initial tokenization phase, BLT removes the need for the heuristic step of converting bytes into a predefined vocabulary of tokens. This conventional step imposes constraints that shape string compression, and is associated with issues such as sensitivity to the specific domain or data modality~\citep{dagangetting}, increased vulnerability to noisy inputs~(\S\ref{sec:robustness}), insufficient representation of orthographic features~\citep{edman2024cute}, and unequal treatment of different languages~\citep{liang2023xlm,petrov2024language,limisiewicz2024myte}.

Tokenization has historically been crucial because training llms directly on raw bytes incurs excessive computational expenses at scale, primarily due to the resulting elongated sequence lengths~\citep{xue2022byt5}. Previous studies have sought to address this issue by introducing more computationally efficient forms of self-attention~\citep{el2020characterbert, clark2022canine} or by adopting attention-free model architectures~\citep{wang2024mambabyte}~(\S\ref{section:related_work}). Nonetheless, such approaches mainly facilitate the training of \textit{small models}. For larger models, most computational resources are expended on the extensive feed-forward network layers operating over every byte, whereas the attention mechanism becomes a less significant contributor to overall cost.

We introduce an adaptive and trainable approach for clustering bytes into \textit{patches}~(\S\ref{section:patching}), accompanied by a novel architecture that integrates both byte-level and patch-level data. In contrast to tokenization, {\textsc BLT} does not impose a predetermined set of patches. Instead, it uses efficient, learnable encoder and decoder modules to project arbitrary byte groupings into latent patch representations. Our experiments demonstrate that this approach leads to a \textit{more} effective use of computational resources compared to models relying on tokenization.

Tokenization-based language models assign identical computational resources to each token, which prioritizes uniformity over efficiency. This approach can be suboptimal because the heuristics used to compress data into tokens do not always correspond to the actual predictive complexity. The foundation of our architecture is the principle that models ought to tailor computational effort according to contextual demands. For instance, predicting the completion of most words requires significantly less computation than generating the opening word of a fresh sentence, as the former typically involves simpler, lower-entropy predictions. This distinction is captured by the design of {\textsc BLT}~(\S\ref{section:architecture}), which integrates three transformer components: two compact, byte-level \textit{local models}, alongside a much larger, global \textit{latent transformer}~(\autoref{fig:arch_high_level}). 

To facilitate dynamic compute allocation and decide how bytes are grouped into patches, {\textsc BLT} organizes the input based on the entropy associated with predicting the next byte. This leads to context-aware groupings in which each segment comprises bytes exhibiting similar levels of information density.

We conduct the first flop-controlled scaling investigation of byte-level models at scales reaching 8B parameters and 4T training bytes, demonstrating the feasibility of end-to-end training directly from bytes, eliminating the need for predefined vocabulary tokenization. In terms of training efficiency controlled for flops,\footnote{We estimate a model's computational requirement by enumerating the total number of Floating Point OPerations (flops).} {\textsc BLT} achieves parity with Llama 3 but requires up to 50\% fewer flops during inference~(\S\ref{section:scaling}). Additionally, our findings indicate that directly modeling raw bytes delivers substantial benefits in capturing rare data patterns. {\textsc BLT} models exhibit greater robustness to noisy input when compared to tokenizer-based approaches and demonstrate superior proficiency at the character level, as seen in evaluations involving orthographic knowledge, phonology, and low-resource machine translation~(\S\ref{sec:robustness}).

Furthermore, with {\textsc BLT} models, we are able to scale both model capacity and patch size concurrently, all within a fixed inference flops constraint. Employing longer patch sizes typically results in computational savings, thereby enabling expansion of the global latent transformer module since it is invoked less frequently. Through a series of scaling experiments that hold inference flops constant~(\autoref{fig:fixed_inference_scaling}), we find that {\textsc BLT} models exhibit notably improved scaling behavior compared to conventional tokenization-based designs.

To conclude, our work provides several key advancements: 1) We present {\textsc BLT}, a byte latent LLM architecture capable of adaptive compute allocation to enhance flop efficiency; 2) We demonstrate that {\textsc BLT} matches Llama 3 in training flop-controlled performance up to the 8B parameter range, with the capability to yield up to 50\% improved flop efficiency when accepting slight evaluation tradeoffs; 3) Our approach with {\textsc BLT} introduces a novel scalability axis for LLMs, allowing the model size to grow while upholding a constant inference computation budget; 4) We show that {\textsc BLT} models exhibit greater robustness to noisy inputs and heightened sensitivity to sub-word data features often overlooked by token-centric LLMs. The implementation for both training and inference with {\textsc BLT} is publicly available at~\url{https://github.com/facebookresearch/blt}.

\section{Patching: From Individual Bytes to Groups of Bytes}
\label{section:patching}

Dividing byte sequences into \textit{patches} enables {\textsc BLT} to assign computational resources adaptively according to input context. Figure~\ref{fig:patching_types} illustrates multiple approaches to grouping bytes into patches. Mathematically, a patching function $f_p$ takes an input byte sequence $\pmb{x}=\{x_i,|i=1,\ldots n\}$ of length $n$ and partitions it into a sequence of $m < n$ patches $\pmb{p}=\{p_j|j=1,\ldots,m\}$, assigning each $x_i$ a value from \{0,1\}, where 1 signals the initiation of a new patch. For both token-centric and patch-centric Transformer models, the main computational burden is governed by the number of times the core Transformer processes data—that is, by how many patches are formed using the specific patching function. Therefore, the expected or average number of bytes per patch (termed the \textit{patch size}) becomes the critical determinant of computational requirements during both model training and inference phases~(\S\ref{section:flops}) when a particular patching function is used.

We subsequently describe three distinct patching strategies: segmenting by a constant byte count per patch~(\S\ref{section:static-patch}), using whitespace-delimited patching~(\S\ref{section:white-patch}), and employing adaptive patching based on byte-wise entropy from a lightweight byte-level language model~(\S\ref{section:dyn-patch}). Additionally, we examine incremental patching and distinguish between the mechanics of tokenization versus patching~(\S\ref{section:bpe-patch}).

\subsection{Strided Patching Every K Bytes}
\label{section:static-patch}

A common and intuitive method for organizing bytes is to segment them into patches of a uniform length $k$, as adopted in MegaByte~\citep{yu2023megabyte}. This approach, using a consistent stride, simplifies both the training and inference processes, and allows for easy adjustment of the mean patch size, which in turn facilitates regulation of computational (flop) requirements. Nevertheless, such a patching approach presents notable drawbacks. One limitation is that computational resources are not adaptively focused on regions where they are most beneficial: for instance, a transformer step $j$ may be squandered when processing trivial content like whitespace in code, or, conversely, too little computation may be allocated to patches densely packed with critical information, as in mathematical expressions. Furthermore, this technique can result in irregular and context-insensitive splitting of identical byte sequences—such as partitioning the same word in inconsistent ways.

\subsection{Space Patching}
\label{section:white-patch}

\citet{slagle2024spacebyte} introduces an enhanced approach to strided patching by generating new patches immediately after encountering any space-like bytes\footnote{
    Space-like bytes are defined as any byte that is not a latin character, digit, or utf-8 continuation byte. In addition, each patch must contain at least one non space-like byte. }, which frequently coincide with natural linguistic boundaries in numerous languages. Within the Space patching technique, a latent transformer step (i.e., a greater computational budget in terms of flops) is specifically allocated to each word. This mechanism guarantees consistency in the patching of words across different sequences and directs additional computational resources toward the more challenging predictive points, typically found just after space characters. For instance, identifying the initial byte of the response in the question ``Who composed the Magic Flute?\uline{\hspace{.6em}}'' poses a greater prediction difficulty compared to subsequent bytes following the initial ``M'', since the first character crucially narrows down the set of likely continuations, rendering the completion ``Mozart'' much easier to anticipate. Nevertheless, the space patching approach is limited in that it does not adapt well to all languages and domains, nor can it dynamically adjust the patch size. In the following, we propose a novel patching strategy inspired by the observation that the first bytes of words are usually hardest to predict, while simultaneously offering an intrinsic means to control patch size.

\subsection{Entropy Patching: Using Next-Byte Entropies from a Small Byte LM}
\label{section:dyn-patch}

Instead of utilizing a rule-based heuristic like whitespace, we employ a data-driven strategy to detect points where next-byte prediction uncertainty is elevated. Specifically, we present \textit{entropy patching}, a method that leverages entropy measurements to determine patch boundaries.

We train a small byte-level auto-regressive language model on the training data for {\textsc BLT} and compute next byte entropies under the LM distribution $p_{e}$ over the byte vocabulary $\mathcal{V}$:

\begin{equation}
H(x_i) = \sum_{v \in \mathcal{V}} p_{e}(x_i=v|\pmb{x}_{<i}) \log p_{e}(x_i=v|\pmb{x}_{<i})
\end{equation}

We investigate two approaches for detecting patch boundaries using entropies $H(x_i)$. The initial method selects points where the entropy exceeds a fixed global threshold, as shown in~\autoref{fig:patching}. The alternative method locates points with entropy values significantly higher than the preceding value. This latter strategy can also be viewed as flagging locations where entropy diverges from an approximately monotonically decreasing trend within the patch.

\begin{align*}
    \text{Global Constraint}\hspace{1cm}H(x_t)& > \theta_g\\
    \text{Approx. Monotonic Constraint}\hspace{1cm}H(x_t)& - H(x_{t-1}) > \theta_r
\end{align*}

During the dataloading process, a streamlined preprocessing phase is employed to determine patch boundaries. Unlike~\citet{nawrot-etal-2023-efficient}, who utilize a trained classifier to predict patch boundaries based on entropy, our approach is more lightweight. In our experimental evaluation~(\S\ref{section:experiments}), we provide a comparative analysis of these two techniques for classifying bytes as either low or high entropy.

\subsection{The Byte-Pair Encoding (BPE) Tokenizer and Incremental Patching}
\label{section:incremental-patching}
\label{section:bpe-patch}

A wide range of current llms, such as our chosen baseline Llama 3, employ subword tokenizers like bpe~\citep{gage1994new, sennrich-etal-2016-neural}. In this context, we define ``tokens'' as collections of bytes taken from a \textit{finite} vocabulary that is established before model training begins. In contrast, ``patches'' denote adaptively created groupings of sequences that do not rely on a predefined vocabulary. The key distinction between tokens and patches lies in the fact that, with tokens, the model lacks explicit access to the raw byte-level features.

An important advancement introduced by {\textsc BLT} compared to tokenization-based architectures lies in its redefinition of the balance between vocabulary size and computational requirements. Conventional large language models tend to face a dilemma: expanding the vocabulary typically results in larger average tokens, which reduces the sequence length needed, but also necessitates a substantial increase in the output dimensionality of the model’s final projection layer. As a result, tokenization-centric methods have limited flexibility to meaningfully adjust token sizes or inference expenses. As an illustration, Llama 3 raises the mean token length from 3.7 to 4.4 bytes, but this improvement comes at the expense of expanding its embedding table by a factor of four relative to Llama 2.

During generation, {\textsc BLT} must determine at each position in the byte sequence whether it has reached a patch boundary, as this choice influences when additional computation via the Latent Transformer is applied. This boundary-detection must occur without reference to the future elements of the sequence, since those elements have not yet been generated. Consequently, the patching process cannot rely on knowledge of subsequent bytes when segmenting the sequence. More precisely, a patching function $f_p$ is said to possess the property of incremental patching if the following holds:
$$f_p(\pmb{x}_{<i}) = f_p(\pmb{x})_{<i}$$
The commonly used bpe segmentation method does not fulfill incremental patching, since it may produce different tokenizations for the same prefix contingent on the tokens that follow. As such, it fails to meet the criterion specified above.\footnote{Introducing a unique delimiter token to mark patch boundaries can render bpe incremental, but this approach comes at the cost of increased byte sequence length.}

\section{{\textsc BLT} Architecture}
\label{section:architecture}
{\textsc BLT} consists of a substantial autoregressive language model at the global level, which processes patch representations. In addition, it utilizes a pair of more compact local models: one responsible for converting byte sequences into patches, and another for reconstructing bytes from the patch representations~(\autoref{fig:arch_high_level}).

\subsection{Latent Global Transformer Model}

\textit{The Latent Global Transformer} refers to an autoregressive transformer, denoted as $\mathcal{G}$, comprising $l_{\mathcal{G}}$ layers. This model transforms a sequence of latent input patch embeddings, $p_j$, into a corresponding sequence of output patch embeddings, $o_j$. In this work, we consistently use $j$ to represent patch indices, while $i$ is reserved for byte indices. The global model leverages a block-causal attention mask~\citep{dubey2024llama}, ensuring that attention is limited to the current patch and all preceding patches within the same document. The computational workload, both during pre-training and inference, is primarily concentrated in this module. Consequently, judicious selection of when to apply the global model enables adaptive allocation of computational resources across various segments of the input and output, modulated by their respective complexities.

\subsection{Local Encoder}
\label{section:local}

The \textit{Local Encoder Model}, represented as $\mathcal{E}$, is designed as an efficient, transformer-based encoder consisting of $l_{\mathcal{E}}$ layers, where $l_{\mathcal{E}}$ is significantly smaller than $l_{\mathcal{G}}$. Its primary function is to rapidly project an input sequence of bytes $b_i$ into informative patch-level representations $p_j$. Notably, unlike conventional transformer structures, this model incorporates a cross-attention layer subsequent to each transformer block to aggregate the byte-level features into patch-level representations (see~\autoref{fig:crossattn}).

Initially, each input byte $b_i$ is mapped into an embedding space via a matrix in $\mathbb{R}^{256 \times h_{\mathcal{E}}}$, yielding embedded vectors $x_i$. Optionally, these embeddings may be enriched by integrating hash-embeddings as described in~\S\ref{sec:hash-emb}. 

The model then applies a succession of alternating transformer layers and cross-attention blocks (\S\ref{sec:enc-cross-attn}) to progressively develop the patch representations $p_i$, which serve as inputs to the global transformer module $\mathcal{G}$. Each transformer layer leverages a \textit{local block causal} attention mask; in this scheme, each byte attends to a predetermined number $w_{\mathcal{E}}$ of previous bytes, potentially spanning variable patch boundaries but strictly constrained within a single document. Subsequent sections provide further details regarding the embedding mechanism and the architecture of the cross-attention module.

\subsubsection{Encoder Hash n-gram Embeddings}
\label{sec:ngram-emb}
\label{sec:hash-emb}
An essential factor in developing resilient and informative representations at each position $i$ involves capturing the context of prior bytes. In {\textsc BLT}, this is implemented by encoding not just the standalone byte $b_i$ but also by embedding it within the context of its surrounding byte n-grams. Specifically, at every position $i$, we begin by generating the corresponding byte-grams

\begin{align}
    g_{i,n}=\{b_{i-n + 1},\ldots, b_i\}
\end{align}

for each byte position $i$ and for $n$ ranging from three to eight.\footnote{ Byte-grams are excluded when their size $n$ exceeds the position $i$ (i.e., if $i<n$). }

Subsequently, we propose hash $n$-gram embeddings, which assign every possible byte $n$-gram to an entry in an embedding table $E_{n}^{hash}$—distinct for each $n$ in the set $\{3,4,5,6,7,8\}$—using a hash function, following~\citep{bai2010learning}. The embeddings obtained from this process are summed with the corresponding byte embedding, after which the sum undergoes normalization and is input into the local encoder. The enriched embedding at position $i$ is computed as follows:

\begin{align}
e_i &= x_i + \sum_{n = 3,...,8}E_{n}^{hash}(\text{Hash}(g_{i,n})) \\
\text{where, Hash}(g_{i,n}) &= \text{RollPolyHash}(g_{i,n}) \% |E_{n}^{hash}|
\end{align}

The value $e_i$ is scaled by dividing it by the total number of $n$-gram sizes incremented by one, and RollPolyHash, as detailed in Appendix~\ref{appendix:rollpolyhash}, is employed in our approach. Section~\ref{section:ablations} presents an analysis of how varying $n$-gram hash embedding parameters, including choices of $n$ and the size of the embedding table, influence trends in the scaling law under controlled flop regimes. Beyond using hash-based $n$-gram embeddings, we also tested frequency-based $n$-gram embeddings; further information regarding these experiments can be found in Appendix~\ref{appendix:ngrams}.

\subsubsection{Encoder Multi-Headed Cross-Attention}
\label{sec:enc-cross-attn}
Our approach largely mirrors the input cross-attention mechanism utilized in the Perceiver architecture~\citep{jaegle2021perceiver}, but a key distinction is that the latent representations now correspond to representations of variable patches rather than being derived from a predetermined set of latents~(\autoref{fig:crossattn}). Additionally, each patch's latent only attends to the bytes comprising that particular patch. Within this module, a query vector is associated with each patch $p_j$; this vector is generated by pooling together the byte representations within $p_j$, and subsequently applying a linear transformation, $\mathcal{E}_{C} \in \mathbb{R}^{h_{\mathcal{E}} \times (h_{\mathcal{E}} \times U_{\mathcal{E}})}$, where $U_{\mathcal{E}}$ designates the count of cross-attention heads in the encoder. More specifically, denoting the bytes that constitute patch $p_j$ by $f_{\text{bytes}}(p_j)$, we then compute

\begin{align}
P_{0,j} &= \mathcal{E}_{C}(f_\text{bytes}((p_j)), f~ \text{is a pooling function} \\
P_l &= P_{l-1} + W_o\left(\text{softmax}\left(\frac{QK^T}{\sqrt{d_k}}\right)V\right) \\
\text{where } Q_j &= W_q(P_{l-1,j}), K_i = W_k(h_{l-1,i}), V_i = W_v(h_{l-1,i}) \\
h_l &= \text{Encoder-Transformer-Layer}_l(h_{l-1})
\end{align}

In this context, $P \in \mathbb{R}^{n_p \times h_{\mathcal{G}}}$ denotes the representations of $n_p$ patches intended for input to the global model, where each patch vector is derived by aggregating its associated byte embeddings $e_i$. The matrices $W_q$, $W_k$, $W_v$, and $W_o$ serve as linear transformations generating queries, keys, values, and outputs, respectively. Here, the keys and values are computed via projections of the byte hidden states $h_i$ from the previous layer, and from $e_i$ in the initial layer. The patch-aware masking approach ensures that each query $Q_j$ focuses attention solely on the keys and values belonging to the corresponding patch $j$. To accommodate multi-headed attention, the patch vectors $P_l$ are kept as several heads of size $h_{\mathcal{E}}$ during cross-attention—recognizing that patch embedding dimensionality $h_{\mathcal{G}}$ typically exceeds $h_{\mathcal{E}}$—and the outputs of these heads are subsequently concatenated to attain dimension $h_{\mathcal{G}}$. Moreover, a pre-LayerNorm operation is applied to the query, key, and value sequences, and no positional encodings are included within this cross-attention mechanism. Lastly, a residual connection is employed encompassing the entire cross-attention layer.

\subsection{Local Decoder} 

Analogous to the design of the local encoder, the local decoder $\mathcal{D}$ also utilizes a transformer architecture with a small number of layers ($l_{\mathcal{D}} << l_{\mathcal{G}}$), making it lightweight. This module is responsible for mapping a sequence of global patch representations $o_j$ back to raw bytes $y_i$. Functioning autoregressively, the local decoder generates each byte based on all previously generated bytes, taking the associated hidden states from the local encoder as input. Its architecture comprises $l_{\mathcal{D}}$ layers, with each layer alternating between a cross-attention mechanism and a transformer block. In this setup, the cross-attention module is applied before the transformer block to align the patch-level representations with byte-level tokens, followed by the transformer layer operating over the constructed sequence of bytes.

\subsubsection{Decoder Multi-headed Cross-Attention} 

In the decoder cross-attention, the roles of the queries and key/values are interchanged i.e. the byte-representations are now the queries, and the patch representations are now the key/values. The initial byte-representations for the cross-attention are initialized as the byte embeddings from the last encoder layer i.e. $h_{l_{\mathcal{E}}}$. The subsequent byte-representations for layer $l$, $d_{l,i}$ are computed as:

\begin{align}
D_0 &= h_{l_{\mathcal{E}}} \\
B_l &= D_{l-1} + W_o\left(\text{softmax}\left(\frac{QK^T}{\sqrt{d_k}}\right)V\right), \\
  &\text{where } Q_i = W_q(d_{l-1,i}), K_i = W_k(\mathcal{D}_C(o_j)), V_i = W_v(\mathcal{D}_C(o_j)) \\
  D_l &= \text{Decoder-Transformer-layer}_l(B_l)
\end{align}

Here, $W_k$ and $W_v$ serve as the projection matrices for keys and values, respectively. These matrices are used within a linear transformation combined with the split operation $\mathcal{D}_C$, which is applied to the final patch representations $o_j$ generated by the global model. The matrix $W_q$ is used to project queries, operating on the byte representations $d_{l-1}$ from the previous layer in the decoder transformer, or $h_{l_{\mathcal{E}}}$ for the initial layer. The output projection matrix is denoted by $W_o$. The resulting matrix $B$ thus resides in $\mathbb{R}^{h_{\mathcal{D}} \times n_b}$, where $n_b$ designates the total number of output bytes. The subsequent representations of the decoder, $D_l$, are obtained by applying a decoder transformer layer to the cross-attention block output $B$. Mirroring the approach in the local encoder cross-attention module, we deploy multiple attention heads, employ pre LayerNorm, omit positional embeddings, and incorporate a residual connection encircling the cross-attention operation.

\section{Experimental Setup}
\label{section:experiments}
To rigorously assess the performance of {\textsc BLT} relative to tokenization-based models, we implement carefully controlled experiments. We ensure that {\textsc BLT} does not benefit from any inadvertent advantages such as longer context lengths.

\subsection{Pre-training Datasets} For all model configurations examined in this study, we utilize two primary pre-training datasets: 1) The Llama 2 dataset~\citep{touvron2023llama}, containing 2 trillion tokens sourced from a diverse selection of publicly available materials, which undergo rigorous cleaning and filtering to ensure higher data quality; and 2) {\textsc BLT}-1T: a new compilation with 1 trillion tokens, also sourced from public domains and including a portion of the pre-training corpus released by Datacomp-LM~\citep{li2024datacomp}. The first dataset serves as the foundation for scaling law analyses to assess the optimal token budget following~\cite{dubey2024llama}, thereby informing the architectural parameters for {\textsc BLT}. The second dataset is employed for full-scale pre-training, facilitating direct comparison with Llama 3 on various downstream benchmarks. Importantly, neither dataset contains data from Meta's products or services. In addition, for our baseline experiments involving tokenizer-based models, we opt for the Llama 3 tokenizer with a vocabulary size of 128K tokens, as it consistently yields better baseline results than the Llama 2 tokenizer within our experimental setting.

\subsection{Entropy Model} For our experiments, the entropy model utilized is a byte-level language model trained on the identical data distribution as the comprehensive {\textsc BLT} model. By default, this model is implemented as a transformer architecture comprising 100M parameters, with 14 layers, a hidden size of 512, and a sliding window attention mechanism spanning 512 bytes. All other hyperparameters mirror those employed in both the local and global transformer baselines. As outlined in \autoref{section:ablations}, we further investigated variations in model scale, architectural design, and receptive field. Notably, for sufficiently constrained receptive fields, it becomes feasible to encode the learned entropy model as a compact lookup table.

\subsection{Entropy Threshold and Equalizing Context Length} For models that employ entropy-driven patching, we determine a patching threshold so as to obtain a specified average \textit{patch size} over the pretraining dataset mixture. In {\textsc BLT}, as opposed to conventional tokenization approaches, the \textit{patch size} is freely adjustable, significantly influencing the model’s effective context window. To ensure fair comparison and to prevent larger patch sizes from yielding disproportionate advantages, we normalize the expected number of bytes per batch. As a result, when employing models with increased patch sizes, we correspondingly shorten the sequence lengths. Specifically, for the Llama 2 dataset, we adopt an 8k byte context, whereas for the {\textsc BLT}-1T dataset, we set the context to an average of 16k bytes, with both settings keeping the average batch size fixed at 16M bytes.

Although the batch size remains fixed, the ratio of bytes to patches fluctuates with the use of dynamic patching approaches during data batching. To enhance efficiency, our {\textsc BLT} training procedure organizes patch batches in a way that circumvents padding within the computationally intensive latent transformer component. This strategy guarantees a uniform patch count per batch. Throughout training, byte sequences are either padded or truncated to 12k bytes for Llama 2 and 24k bytes for {\textsc BLT}-1T datasets, thereby preventing memory overload that can result from exceptionally large patch sequences.

\subsection{Entropy Model Context}
\label{section:entropy-context}
Our empirical investigations indicate that entropy patching leads to increasingly larger patches when applied to structured datasets, such as those involving multiple choice formats (see the patching example for MMLU in \autoref{fig:mmlu_patching}), which tend to display high levels of repetitiveness. This trend can be attributed to the lower entropy associated with repeated elements within the entropy model context. Consequently, for the comprehensive {\textsc BLT}-Entropy experiment utilizing a patch size of 4.5, we refresh the entropy context by inserting new lines and employ an approximate monotonicity constraint, as this approach is less susceptible to the “entropy drift” induced by fluctuations in context length. Notably, this adjustment only influences the method by which we calculate entropy, while our approach for determining the entropy threshold value remains unchanged.

\subsection{FLOPs Estimation} 
\label{section:flops}

Our calculation of transformer flops predominantly adheres to the methodology outlined in Chinchilla~\citep{hoffmann2022training}, which aggregates flops across feed-forward components, qkvo projections within self-attention, and the attention plus output projection steps. However, in contrast to the original procedure, we presume the input embedding stage is realized via an efficient lookup mechanism rather than by dense matrix multiplication, effectively rendering its flop contribution negligible. Consistent with earlier studies, we further assume that the backward pass demands double the number of flops required for the forward pass.

To compute flops \textit{per byte} for {\textsc BLT} models, we add up the flops for the local encoder transformer, the global latent transformer, and the local decoder transformer, together with the cross attention blocks in the encoder and the decoder:

\begin{align}
    \text{FL}_{\text{{\textsc BLT}}} &= \text{Transf. FL}(h_{\mathcal{G}}, l_{\mathcal{G}}, m=n_{ctx}/n_p, V=0) / n_p\\
   &+ \text{Transf. FL}(h_{\mathcal{E}}, l_{\mathcal{E}}, m=w_{\mathcal{E}}, V=0) \\
   &+ \text{Transf. FL}(h_{\mathcal{D}}, l_{\mathcal{D}}, m=w_{\mathcal{D}}, V=256) \\ 
    &+ \text{Cross Attn. FL}(h_{\mathcal{E}}, l_{\mathcal{E}}, m=n_p, r=n_p/k) \cdot k/n_p \\ 
   &+ \text{Cross Attn. FL}(h_{\mathcal{D}}, l_{\mathcal{D}}, m=k, r=k/n_p) 
\end{align}

Here, $n_{ctx}$ denotes the byte-level sequence length, $n_p$ specifies the patch size, $r$ represents the query-to-key/value ratio, and $k$ is defined as the proportion between the patch dimension and byte dimension—this determines how many local model splits are concatenated to assemble the global model representation ($k=2$ as illustrated in \autoref{fig:crossattn}). $V$ signifies the output projection’s vocabulary size, relevant solely for the local decoder. The effective context length, $m$, used for attention calculations, varies based on whether the module processes a byte or patch sequence. Attention flop calculations are adapted for each architectural component based on this context length. Detailed formulations for computing Transformer-FLOPs and Cross-Attention FLOPs can be found in Appendix~\ref{section:flops-appendix}.

\subsection{Bits-Per-Byte Estimation} Perplexity only makes sense in the context of a fixed tokenizer as it is a measure of the uncertainty for each token. When comparing byte and token-level models, following previous work~\citep{xue2022byt5, yu2023megabyte, wang2024mambabyte}, we instead report Bits-Per-Byte (BPB), a tokenizer independent version of perplexity. Specifically:

\begin{align}
    \text{BPB}(x)&= \frac{\mathcal{L}_{CE}(\pmb{x})}{\ln(2)\cdot n_{\text{bytes}}}
\end{align}

Here, the aggregate cross-entropy loss assessing the uncertainty in $\pmb{x}$ is divided by both the number of bytes comprising $\pmb{x}$ and a fixed constant for normalization.

\subsection{Transformer Architecture Hyperparameters} Throughout all transformer components in {\textsc BLT}, including both the local and global models, our architectural choices closely align with those found in Llama~3~\citep{dubey2024llama}. Specifically, in the feed-forward networks, we employ the SwiGLU activation function~\citep{swiglu}, while self-attention modules are equipped with rotary positional embeddings (RoPE)~\citep{RoPe}, where we set $\theta=500000$ as recommended by~\citep{xiong2024effective}. Layer normalization is performed using RMSNorm~\citep{RMSNorm}. For self-attention operations that adopt static, standard masking—such as \textit{block causal} or \textit{fixed-window block causal}—we utilize the Flash attention mechanism~\citep{dao2022flashattention}, and restrict the window size to 512 for masks of fixed width. To address the dynamically changing, patch-sensitive masks present in our cross-attention layers, we adopt Flex Attention\footnote{\url{https://pytorch.org/blog/flexattention}}, which enables efficient fused kernel computations and leads to a substantial acceleration of training.

\subsection{{\textsc BLT}-Specific Hyperparameters} 

To evaluate the performance of {\textsc BLT} models, we design experiments along two main axes: analysis of scaling behavior and assessment on downstream tasks. Our experiments encompass model sizes of 400M, 1B, 2B, 4B, and 8B parameters, with architectural specifics provided in Appendix~Table~\ref{tab:arch}. 

For the initialization of queries in the first cross-attention layer of the local encoder, we utilize max-pooling. Across all {\textsc BLT} variants, we implement $500,000$ hashes using a single hash function, with n-gram lengths varying between 3 and 8. The learning rate is fixed at $4e-4$ for every model. This setting aligns the learning rate between token models and {\textsc BLT}, based on hyperparameter sweeps conducted at the 400M and 1B scales within the interval $1e-3$ to $1e-4$, which indicated this shared rate yielded optimal outcomes. For the scaling law experiments on Llama-2 data, batch sizes follow the recommendations of~\cite{dubey2024llama}, adjusted to byte-equivalent batch sizes where necessary. 

Regarding optimization, we employ the AdamW optimizer~\citep{adamw}, with hyperparameters $\beta_1=0.9$, $\beta_2=0.95$, and $\epsilon=10^{-8}$. Training utilizes a linear warm-up over 2000 steps, followed by a cosine learning rate decay to zero. We set the weight decay coefficient at 0.1, and apply global gradient clipping capped at 1.0.

\section{Scaling Trends}
\label{section:scaling}

We provide a comprehensive overview of the scaling behavior exhibited by byte-level models, offering insights pertinent to the ongoing development and scaling of {\textsc BLT} architectures. Our scaling investigation addresses gaps in existing research on byte-level modeling through the following contributions: (a) conducting a comparison of scaling dynamics within the regime where computational efficiency is optimized; (b) training 8B-parameter models on large-scale datasets (reaching up to 1T tokens/4T bytes) and assessing their performance on various downstream tasks; and (c) analyzing scaling behavior under constant inference cost constraints. Subsequent sections will explore in detail the unique benefits conferred by directly modeling byte-sequences.

\subsection{Parameter Matched Compute Optimal Scaling Trends}
\label{section:parameter-matched-scaling}
We employ the Llama 2 dataset to train both \textit{compute-optimal} bpe and {\textsc BLT} models, each spanning four distinct model sizes from 1B to 8B parameters. Subsequently, we chart the relationship between the training flops expended and the language modeling outcomes over a representative portion of the training data mixture. The bpe models are optimized with respect to the parameter-to-data ratio, as prescribed by Llama~3~\citep{dubey2024llama}, thereby aligning with the \textit{compute-optimal} paradigm. This configuration is theoretically poised to maximize performance for a given training compute allocation~\citep{hoffmann2022training}, thus serving as a strong reference point for our model comparisons. Each bpe model has an associated {\textsc BLT} variant, trained on identical data and featuring a Latent Transformer architecture precisely mirroring that of the bpe model counterpart.

As depicted in~\autoref{fig:scaling-compute-opt} (right), {\textsc BLT} models consistently equal or surpass the performance of bpe models, a pattern that persists even as we increase model capacity and computational resources. To our knowledge, {\textsc BLT} represents the first byte-level Transformer design to attain comparable scaling properties to those of BPE-based models within compute-optimal conditions. This finding supports our hypothesis that the ideal balance between model parameters and training compute identified for bpe models is also applicable to {\textsc BLT}, or at least approximates it closely.

Both architectural enhancements and dynamic patching play a vital role in achieving bpe scaling behaviors. In ~\autoref{fig:scaling-compute-opt} (left), we present a comparison between models utilizing space-patching techniques and Llama 3. For this, we estimate the performance of SpaceByte \citep{slagle2024spacebyte} by implementing {\textsc BLT} space-patching while excluding n-gram embeddings and cross-attention mechanisms. While SpaceByte demonstrates an improvement over Megabyte, there is still a significant gap compared to Llama 3. \autoref{fig:scaling-compute-opt} (right) demonstrates the contribution of both design modifications and the integration of dynamic patching. At scale, {\textsc BLT} models achieve results comparable to leading tokenizer-based systems such as Llama 3.

Furthermore, our results indicate that for models reliant on tokenization, selecting the Llama-3 tokenizer leads to better performance compared to employing the Llama-2 tokenizer, even when both are trained on identical datasets.

In summary, our {\textsc BLT} framework demonstrates characteristics intermediate between those of Llama 2 and Llama 3, particularly when much larger patch sizes are utilized. While the average token sizes for the bpe tokenizers of Llama 2 and 3 are 3.7 and 4.4 bytes respectively, {\textsc BLT} achieves comparable scaling behaviors with mean patch sizes of 6 and even up to 8 bytes. Since inference flops scale inversely with the average patch size, selecting a patch size of 8 bytes can result in nearly a 50\% reduction in inference computation. Furthermore, increasing patch sizes appears to enhance performance as both the scale of the model and the dataset are enlarged. Notably, {\textsc BLT} with a patch size of 8 begins with a considerably lower performance than the bpe Llama 2 at a 1B parameter scale but ultimately surpasses the bpe baseline at 7B parameters. These findings imply that adopting larger patch sizes could yield even better outcomes as both model scale and training resources increase, suggesting that still greater patch sizes might be viable with further scaling.

\subsection{Beyond Compute Optimal Task Evaluations}
\label{section:evals}
To further investigate model scaling behavior, we train an 8B {\textsc BLT} model in a regime exceeding the compute optimal ratio, utilizing the more extensive and higher-quality {\textsc BLT}-1T dataset. The model’s abilities are then evaluated across a variety of widely recognized classification and generation benchmarks. For benchmarking, we focus on the following areas: common sense reasoning, world knowledge, and code generation.
\paragraph{Classification tasks} encompass ARC-Easy (0-shot)~\citep{Eval_ARC}, Arc-Challenge (0-shot)~\citep{Eval_ARC}, HellaSwag (0-shot)~\citep{Eval_hellaswag}, PIQA (0-shot)~\citep{Eval_piqa}, and MMLU (5-shot)~\citep{hendrycks2020measuring}. For assessment, we use a prompt scoring approach based on computing the likelihoods for each answer choice, and accuracy is averaged and reported.
\paragraph{Coding related generation tasks:} We assess the model’s capacity for Python code generation by presenting pass@1 results on MBPP (3-shot)~\citep{Eval_MBPP} and HumanEval (0-shot)~\citep{Eval_HumanEval}.

In \autoref{tab:evals}, we present a comparison among three models trained on the {\textsc BLT}-1T dataset: a model using the bpe Llama 3 tokenizer,\footnote{We opt for the Llama 3 tokenizer, featuring a 128k token vocabulary, because it yields superior performance relative to the 32k vocabulary of Llama 2.} along with two variants of the {\textsc BLT} model. The first variant applies a space-patching strategy ({\textsc BLT}-Space), while the second adopts an entropy-driven patching strategy ({\textsc BLT}-Entropy), incorporating an approximate monotonicity constraint and resetting the entropy model's context at newlines (following the approach outlined in~\autoref{section:entropy-context}). Training is conducted for all three models using an equivalent computational budget measured in flops. Notably, for {\textsc BLT}-Entropy, we further modify the inference procedure by lowering the entropy threshold from 0.6 to 0.1, a change that boosts task performance, albeit requiring more inference steps.

The {\textsc BLT}-Entropy model achieves superior results compared to the Llama 3 model in 4 out of 7 evaluated tasks, despite both being trained with an equivalent amount of data. This performance gain can likely be attributed to two factors: (1) enhanced utilization of training compute facilitated by dynamic patching, and (2) the explicit representation of information at the byte level rather than the token level.

Conversely, {\textsc BLT}-Space demonstrates lower performance than the Llama 3 tokenizer across all but one task. However, its average patch size of 6 bytes leads to a notable decrease in inference flops. By contrast, both bpe and entropy-patching models yield similar average patch sizes, around 4.5 bytes on the mixed training dataset. Given an equivalent training budget, the model with the larger patch size processes 30\% more data than the others. This increased data coverage could potentially move {\textsc BLT} further from the compute-optimal regime.

\subsection{Patches Scale Better Than Tokens} The {\textsc BLT} models enable joint scaling of both the model architecture and the patch size without increasing the overall computation required for training or inference, and without needing more training data. This flexibility in enlarging patch size distinguishes patch-based approaches from traditional fixed-vocabulary token models, which face inherent efficiency limitations, as outlined in Section~\ref{section:bpe-patch}. By processing longer patches, computational demand decreases, freeing resources that can then be used to expand the global latent transformer, since it needs to execute fewer times.

To evaluate whether larger models executing fewer steps on bigger patches outperform smaller models with more steps, we perform a fixed inference scaling analysis. Using models with 400 million and 3.6 billion parameters equipped with the Llama 2 tokenizer as baselines, we identify flop-matched counterparts utilizing the Llama 3 tokenizer and {\textsc BLT}-Entropy models configured with mean patch sizes of 6 and 8 bytes on the same training datamix (model configurations detailed in~\autoref{tab:fixed-inf-params}). For models employing a patch size of 8, we implement 3 encoder layers instead of a single layer. Each model is then trained under multiple training flop constraints.

As depicted in \autoref{fig:fixed_inference_scaling}, {\textsc BLT} models demonstrate superior scaling properties relative to tokenization-based systems across both types of inference flop limits. While BPE models initially excel with modest training resources, they are soon overtaken by {\textsc BLT} models just beyond the point of compute-optimal training. In practical settings, investing additional computation during the pretraining phase can yield a model that delivers stronger performance under a set inference constraint. A notable case is the 8B parameter class, such as Llama 3.1, which underwent pretraining on data volumes roughly two orders of magnitude greater than those typically considered optimal for its scale.

The point at which {\textsc BLT} surpasses token-based models now occurs closer to the compute-optimal regime in the case of higher flop class models, shifting from three times to approximately 2.5 times the compute-optimal budget. The model using patch size 8 within this higher flop category also demonstrates a more pronounced scaling behavior, overtaking other configurations at an earlier stage. As detailed in Section~\ref{section:parameter-matched-scaling}, models with increased patch sizes tend to align more closely in performance with BPE models as the overall model scale grows. We partly attribute this phenomenon to the reduction in the proportion of total flops allocated to the byte-level Encoder and Decoder modules, as these components exhibit a slower scaling rate compared to the Latent Transformer. When expanding the total parameter count by a factor of 20—from 400M to 8B—the increase in local model parameters for {\textsc BLT} is only about two-fold. This distinction matters, as increases in patch size impact flops within the patch Latent Transformer but do not affect the byte-level modules. Consequently, when advancing to a larger scale, the model size for {\textsc BLT}-Entropy with ps=8 rose only marginally, from 1.6x to 1.7x relative to the Llama 2 model.

To summarize, our analysis of scaling with respect to patch length shows that the {\textsc BLT} patch-based model benefits from joint increases in both the size of the patches and the overall model. These advantageous scaling patterns not only persist but may become even more pronounced as models grow larger.

\section{Byte Modeling Improves Robustness}
We further assess the robustness of {\textsc BLT} in relation to token-based architectures that do not incorporate byte-level signals, and introduce a strategy for converting pretrained token-based models to operate at the byte level.
\label{sec:robustness}

\subsection{Character-Level Tasks}

One of the initial incentives for developing models that operate at the byte level was their inherent resilience to noise present at the byte layer in the input, in addition to their ability to recognize and make use of the internal structure of tokens—a capability with which existing tokenizer-based approaches often have difficulty. To assess these aspects, we conduct supplementary experiments using benchmarks designed to test both robustness against input disturbances and sensitivity to character-level details, encompassing not just English but also a range of languages, and including characters such as digits and phonetic symbols. These experimental outcomes are summarized in Table~\ref{tab:char_tasks}.

\paragraph{Noisy Data} To assess model robustness, we generate altered versions of the standard classification datasets outlined in Section~\ref{section:evals}, enabling direct comparison between the tokenizer-based models and {\textsc BLT}. For this purpose, we apply five separate character-level noise transformations: (a) \textit{AntSpeak}: Renders the complete text in uppercase letters, with characters separated by spaces. (b) \textit{Drop}: Randomly deletes 10\% of all characters present in the text. (c) \textit{RandomCase}: Randomly assigns uppercase to half the characters and lowercase to the remainder, distributing case assignments throughout the text. (d) \textit{Repeat}: Selects 20\% of the characters at random and repeats them up to four consecutive times. (e) \textit{UpperCase}: Converts the entire textual input to uppercase letters. Each noising technique is applied individually during evaluation, either to the prompt, the completion, or both, and results are averaged across these scenarios. In Table \ref{tab:char_tasks}, our findings on noised HellaSwag~\citep{Eval_hellaswag} show that {\textsc BLT} consistently exhibits greater robustness than its tokenizer-based counterparts, demonstrating a mean improvement of 8 points relative to a model with identical training data, and achieving higher scores even when compared with the Llama 3.1 model, despite the latter being trained on a substantially larger corpus.

\paragraph{Phonology - Grapheme-to-Phoneme (G2P)} We evaluate the proficiency of {\textsc BLT} in converting a sequence of graphemes (alphabetic symbols that represent words) into their corresponding phoneme transcriptions, which reflect the word's pronunciation. Table \ref{tab:char_tasks} shows the 5-shot G2P task results using the Phonology Bench~\citep{suvarna2024phonologybench}. The data indicate that {\textsc BLT} achieves better performance compared to the baseline model based on the Llama 3 1T tokenizer for this specific task.

\paragraph{CUTE}
To evaluate character-level comprehension, we test {\textsc BLT} on the CUTE benchmark~\citep{edman2024cute}. This suite consists of multiple tasks, which can be grouped into three primary types: tasks that examine compositional understanding, those probing orthographic similarity, and those that involve sequence manipulation. CUTE presents considerable difficulty for most models reliant on tokenization, as these models generally retain some awareness of token orthography but often fail to leverage this capacity for concrete text manipulation. According to Table~\ref{tab:char_tasks}, {\textsc BLT}-Entropy surpasses both Llama 3 models based on BPE by over 25 points on this benchmark. Notably, {\textsc BLT} excels at tasks centered on character manipulation, attaining a score of 99.9\% on both spelling tasks. This pronounced performance boost is observed despite {\textsc BLT} being trained on data that is 16 times smaller compared to Llama 3.1, which suggests that BPE-based approaches encounter fundamental obstacles in acquiring robust character-level representations. Figure \ref{fig:cute_examples} showcases instances in which the Llama 3 tokenizer-based model is inadequate, while {\textsc BLT} handles the tasks successfully. The only tasks where BPE-based models outperform ours are word deletion and insertion. Although a byte-level model may find such word-based manipulations less straightforward, the difference in performance is not substantial, and it could be more tractable to scale up from character-level to word-level operations than the reverse. The evaluation procedure and prompts strictly follow those provided by Huggingface for all experiments. Additional prompt engineering may enhance the performance of BPE-based models.

\paragraph{Low Resource Machine Translation} We assess the performance of {\textsc BLT} on translation tasks involving both directions—into and out of—six major language families and twenty-one low-resource languages employing diverse scripts, as sourced from the FLORES-101 benchmark~\citep{goyal2022flores}. SentencePiece BLEU results are summarized in Table~\ref{tab:flores}. The findings indicate that {\textsc BLT} yields stronger results than a counterpart trained with the Llama 3 tokenizer, presenting an overall improvement of 2 BLEU points for translations into English and a 0.5-point increase for translations from English. While performance on widely used language pairs is either on par with or marginally superior to that of Llama 3, {\textsc BLT} notably surpasses Llama 3 on various pairs within underrepresented language families. These outcomes highlight the advantage of byte-based modeling in capturing the generalization potential required for rare and complex byte sequences.

\subsection{Training {\textsc BLT} from Llama 3}
\label{sec:distillation}
We investigate an approach where {\textsc BLT} models benefit from pre-existing, pretrained tokenizer-based architectures to facilitate more efficient and rapid training. This is accomplished by initializing the global transformer parameters within {\textsc BLT} using those from a pretrained Llama 3.1 model. During subsequent training, the global transformer parameters are updated using a learning rate that is one-tenth that applied to both the local encoder and local decoder models, maintaining the same number of optimization steps used in the original Llama 3 training. Table~\ref{tab:distill} contrasts this method with a standard {\textsc BLT} baseline. Results demonstrate that initializing {\textsc BLT} with Llama 3.1 leads to clear improvements over both the Llama 3 baseline and the baseline {\textsc BLT}, each of which was trained using an equivalent number of compute resources. Furthermore, a comparison with the {\textsc BLT}-Entropy variant from Table~\ref{tab:evals}—which was trained on a much larger dataset consisting of 1T tokens—shows that {\textsc BLT} initialized from Llama 3.1 still surpasses it on the MMLU evaluation, indicating that this strategy can substantially reduce training flops while maintaining strong performance.

This approach can be interpreted as converting models that utilize tokenizers into ones that do not, thereby adapting a pre-trained LLaMA 3.1 model into a {\textsc BLT} framework. For a thorough comparison, Table \ref{tab:distill} presents results for both the original LLaMA 3.1 model, trained on 15T tokens, and the {\textsc BLT} model derived from LLaMA 3. Although our approach yields only marginally lower results on MMLU and HumanEval, more substantial performance reductions are observed on the remaining benchmarks. These outcomes indicate that there is room for further exploration to make optimal use of the pre-trained model, with particular emphasis on enhancing data mixture strategies and fine-tuning other hyperparameters.

\section{Ablations and Discussion}
\label{section:ablations}
This section presents ablation studies that rationalize the architectural decisions behind {\textsc BLT}, as well as the selected patching method and hyper-parameter settings for the 8B parameter {\textsc BLT} model trained using the {\textsc BLT}-1T dataset.

\paragraph{Entropy Model Hyper-parameters} To investigate how altering entropy model capacity and context window size influences scaling behavior, we train entropy transformer models at the byte level, spanning model sizes from 1m up to 100m parameters and utilizing a range of context window lengths from 64 to 512. We visualize the relationship between bpb and training flop scaling through law curves, generated from our $400m$ and $1b$ {\textsc BLT} models, both of which are trained on the Llama-2 dataset; these results are displayed in \autoref{fig:entropy_ablation}. Our observations indicate that increasing either model size or context length generally enhances scaling performance, although the marginal benefits decline for model sizes exceeding 50m parameters.

\paragraph{Types of Patching}

We conduct an ablation study on the four distinct patching strategies outlined in Section~\ref{section:patching}: 1) Strided Patching with strides of 4 and 6, 2) Whitespace-based Patching, 3) Patching utilizing the BPE Tokenizer aligned with the Llama 3 tokenizer, and 4) Entropy-based Patching leveraging a compact byte-level LLM.

Although dynamic patching leads to shorter sequences, we standardize sequence length across all patching strategies to ensure that the context length remains comparable. Consequently, each model processes an equivalent number of bytes per sequence, both during training and inference, minimizing potential confounding effects related to context size. The outcomes of these control experiments are illustrated in Figure~\ref{fig:scaling-compute-opt}. Every patching scheme, aside from static patching, achieves superior performance, with space patching performing nearly as well as dynamic entropy based patching.

\autoref{tab:evals_ablation} provides benchmark results for various {\textsc BLT} model configurations, comparing standard tokenizer-based approaches, space patching, and entropy-based patching. All models were trained on the Llama 2 dataset, with training extended to an ideal number of iterations~\citep{dubey2024llama}. While space patching offers a more straightforward approach—since it does not require deploying an entropy model dynamically during training—the improvements seen with entropy-based patching on scaling experiments (see Section~\ref{section:scaling}) are also evident in downstream benchmark tasks.\footnote{Earlier space patching results were obtained without cross-attention, yet comparable trends hold when cross-attention is introduced.}

\paragraph{Cross-Attention} 
In~\autoref{tab:crossattn_ablations_1b}, we systematically examine the impact of introducing cross-attention at different locations within both the encoder and decoder components of {\textsc BLT}. For encoder cross-attention, we evaluate three strategies for initializing the queries: 1) assigning an identical learned embedding to all global states, 2) utilizing a hash embedding derived from the patch bytes, and 3) applying pooling to the encoder’s hidden states of the patch bytes at the corresponding encoder layer.

Our results indicate that incorporating cross-attention within the \textit{decoder} yields the greatest benefits. While applying cross-attention in the encoder leads to marginal gains, this advantage is observed only when queries are initialized through pooling. Furthermore, our analysis shows that cross-attention is especially advantageous for Common-Crawl data, and its utility increases further when employing larger patch sizes.

\paragraph{n-gram Hash Embeddings} 

We conduct an ablation study using n-gram hash embedding vocabulary sizes of 0, 100K, 200K, and 400K, and the corresponding outcomes are shown in Table~\ref{tab:ngram_ablations_1b}. Our analysis reveals that incorporating hash embeddings enhances results across all evaluated domains, with a particularly notable impact on Wikipedia and Github (yielding a 0.04 bpb improvement versus a 0.01 bpb improvement after 15k training steps at the 8B scale). Reducing the number of hashes from 500K to 300K at 8B led to a negligible performance shift of approximately 0.001 bpb at 15k steps. These findings suggest that hash embeddings play an essential role in enabling {\textsc BLT} to achieve competitive performance with tokenizer-based models; nevertheless, additional increases beyond 300K hashes offer minimal incremental benefit. Furthermore, the enhancements provided by hash embeddings appear largely orthogonal to those from cross-attention mechanisms, as improvements manifest across distinct datasets.

\paragraph{Local Model Hyperparamaters} Table \ref{tab:local_layers} presents an ablation study on different configurations for the number of layers in both the local encoder and decoder. Our findings indicate that, when utilizing hash n-gram embeddings, {\textsc BLT} achieves strong performance with an ultra-lightweight encoder—specifically, one layer—while a more complex, multi-layered decoder is advantageous.

\section{Related Work}
\label{section:related_work}

\textbf{Character-Level RNNs:} Character-level language modeling has long been established as a prominent area of research in neural modeling~\citep{sutskever2011generating,mikolov2012subword,graves2013generating}, primarily because of its natural capability to represent out-of-vocabulary terms without relying on explicit back-off techniques. The work by \cite{kim2016character} introduced models that operate solely on characters at the input stage, employing convolutional layers alongside highway networks, which subsequently supply features to LSTM-based RNNs. This architecture achieved performance on par with the leading RNN-based language models of its time for English, and notably surpassed them on languages with complex morphology, thereby highlighting another critical strength of character-level LLMs. Furthermore, \cite{kenter2018byte} investigated machine comprehension via byte-level LSTM architectures, obtaining superior results over word-based systems on morphologically intricate languages such as Turkish and Russian. Similarly, \cite{zhang2015character} leveraged convolutional models at the character level for classification tasks and demonstrated notable improvements over word-level counterparts in some scenarios. Building on these approaches, \citet{chung2019hierarchical} devised hierarchical LSTM frameworks that incorporate boundary-detection at multiple granularities to expose latent text structures, further enhancing character-level language modeling performance. In contrast to attention-driven architectures, ByteNet~\citep{kalchbrenner2016neural} adopts convolutional layers over characters specifically for machine translation tasks.

\textbf{Character-Level Transformers:} The introduction of transformer architectures leveraging attention mechanisms~\citep{vaswani2017attention}, combined with subword tokenization~\citep{sennrich-etal-2016-neural}, led to notable advancements in the capabilities of neural language models across various benchmark tasks. Despite these improvements, the reliance on word or sub-word units introduces a predetermined level of linguistic abstraction into the models. Seeking to merge the strengths of transformer architectures with the initial successes in character-based language modeling, \citet{al2019character} applied extremely deep transformer networks and incorporated auxiliary objectives to develop character-level transformer models that surpassed the performance of previous LSTM-based character language models. Nonetheless, these character-level transformers continued to lag substantially behind their word-level counterparts. Similarly, GPT-2~\citep{radfordgpt2} found that, even when trained on extensive datasets like the 1 Billion Word Benchmark, byte-level language models were outperformed by models operating at the word level.

Although \cite{Choe2019BridgingTG} established that byte-level LLMs employing transformer architectures can surpass subword-level LLMs with comparable parameter counts, they require substantially greater computational resources and prolonged training times. Likewise, \cite{el2020characterbert} introduced CharFormer, a BERT variant that constructs word representations by performing convolutional operations over character embeddings, and showcased performance gains in the medical domain; however, these gains come at a significant computational cost. In their work, \cite{clark2022canine} proposed CANINE, a 150M parameter encoder-only architecture that processes raw character sequences. At the core of CANINE is a deep transformer stack akin to our global model, supplemented by a combination of a local transformer and strided convolutions for effective downsampling of character-level input. CANINE demonstrated superior results to its token-based encoder-only counterpart (mBERT) on downstream multilingual tasks. The ByT5 model~\citep{xue2022byt5} further examined strategies for byte-level encoder-decoder frameworks, explicitly foregoing patching operations. While ByT5 improved resilience to noisy input and achieved comparable outcomes to token-based models with only a quarter of the data, the absence of patching necessitated computationally expensive attention across all bytes, greatly increasing resource requirements. Representing input at the byte level—as opposed to subword units—inevitably increases sequence length, presenting major challenges to the efficient scaling of such models. Most recently, leveraging the Mamba Architecture~\citep{gu2023mamba}, which sustains a constant-size memory state across extended contexts, \cite{wang2024mambabyte} successfully trained a byte-level Mamba-based model (again omitting patching) and achieved better performance than comparable byte-level transformers under matched FLOP conditions, specifically at the 350M parameter scale, as evaluated by bits-per-byte across multiple datasets.

\textbf{Patching-based approaches:} Utilizing patching methods has been shown to reduce the excessive computational costs, measured in flops, commonly associated with byte-level LLMs while still maintaining their effectiveness. Several studies have reported promising early results with modestly sized models and limited data scales. For instance, \cite{nawrot-etal-2022-hierarchical} introduce static patching mechanisms for downsampling and upsampling, leading to the development of the hourglass transformer, which surpasses other byte-level approaches at the 150M parameter level. Building on this, \cite{nawrot-etal-2023-efficient} present improvements through dynamic patching strategies that include an end-to-end trainable boundary predictor, a boundary predictor informed by particular tokenizers, and an entropy-guided patching strategy akin to {\textsc BLT}. Their findings indicate these methods outperform contemporaneous vanilla transformer architectures on language modeling benchmarks using roughly 40M parameters and 400M tokens. Additionally, \cite{lester2024training} explore training models on sequences compacted via arithmetic coding to achieve higher compression rates than BPE, and by implementing an equal-info windows approach, manage to exceed byte-level model baselines in language modeling, though their method still falls short compared to subword model baselines.

Our research is primarily motivated by, and closely aligns with, the approach in MegaByte~\citep{yu2023megabyte}, a decoder-only, causal LLM that applies a fixed, non-adaptive patching strategy. Specifically, MegaByte concatenates byte-level representations into patches, and later decodes these patches back to bytes using a localized decoding model. The original work demonstrates that MegaByte achieves competitive performance with tokenizer-based architectures at the 1B parameter scale when evaluated on datasets of approximately 400B bytes. In our experiments, we include MegaByte as a baseline and observe that fixed patching methods underperform relative to state-of-the-art, compute-efficient tokenizer-based models under flop-constrained scenarios; we then show how {\textsc BLT} effectively addresses this disparity. Likewise, \cite{slagle2024spacebyte} independently report similar findings regarding MegaByte, proposing enhancements such as extending static patching to include segmentation on whitespace or other space-equivalent byte values, and introducing a localized encoder. They note performance gains over token-based transformers in certain domains (e.g., Github, arXiv) within compute-controlled experiments at the 1B parameter scale. Our evaluation of this approach confirms their results, and additionally, we argue that further innovations in architecture are required for byte-level models to reach parity with advanced token-based models like Llama 3 at larger scales.

\section{Limitations and Future Work}

In this study, our selection of architectural parameters relies on training models for the optimal number of steps specified in Llama~3~\citep{dubey2024llama}. It is important to note, however, that these scaling laws were originally established for BPE-level transformers. As such, applying them to {\textsc BLT} might result in less-than-ideal relationships between data and parameter sizes. Determining appropriate scaling laws tailored to {\textsc BLT}—a direction which could further enhance the efficiency of our architecture—is left as a subject for future investigation. Moreover, our experiments largely involve models containing up to 1B parameters. It remains plausible that as model size increases to 8B parameters and above, the optimal design decisions may shift, potentially enabling superior outcomes at larger model scales.

Current transformer libraries and codebases are optimized primarily for tokenizer-based transformer models. Although our work includes theoretical experiments with matched FLOPs and leverages some optimized modules like FlexAttention to accommodate non-standard transformer layers, the practical runtime of our implementations may still fall short of that achieved by tokenizer-based models. Further enhancements and optimizations could potentially narrow this gap in wall-clock efficiency.

Although {\textsc BLT} employs an independently trained entropy model for patching, integrating the learning of the patching model within a fully end-to-end training framework represents a promising avenue for subsequent research. As detailed in Section \ref{sec:distillation}, we provide preliminary empirical results that suggest the feasibility of converting tokenizer-based models, such as Llama 3, which are pretrained on over 10T tokens, into a byte-oriented format. This is achieved by initializing and freezing the global transformer using their existing parameters. Expanding on this line of inquiry could yield approaches that preserve the advantages of bytefying while potentially enhancing the capabilities of these tokenizer-based models, without necessitating a complete retraining.

\section{Conclusion}
\label{section:conclusion}

In this work, we introduce the Byte Latent Transformer (\textbf{BLT}), an innovative model architecture that challenges the standard reliance on predetermined vocabulary tokenization in large language models. Through the adoption of a dynamic and trainable mechanism for assembling bytes into flexible patches, {\textsc BLT} allocates computational effort in accordance with the underlying data’s intricacy, which yields notable advances in computational efficiency and model robustness. Our comprehensive scaling experiments show that {\textsc BLT} achieves parity with established tokenization-dependent models such as Llama 3 at model sizes up to 8B parameters and on datasets up to 4T bytes, while also being able to substantially decrease inference FLOPs—sometimes by as much as 50\%—at the cost of marginal declines in performance metrics. Additionally, {\textsc BLT} enables a novel approach to model scaling: both model size and patch size can be increased simultaneously while keeping inference costs constant, a feature that offers practical benefits in many real-world compute scenarios. By processing raw byte sequences directly, {\textsc BLT} also enhances robustness to noisy or atypical data and enriches the model’s capabilities to interpret sub-word phenomena. Collectively, these advances establish {\textsc BLT} as a compelling candidate to replace conventional tokenization frameworks, offering a scalable, resilient, and efficient architecture for future adaptable language models.

\end{document}